% 詹姆斯·格雷果里（综述）
% license CCBYSA3
% type Wiki

本文根据 CC-BY-SA 协议转载翻译自维基百科\href{https://en.wikipedia.org/wiki/James_Gregory_(mathematician)}{相关文章}

詹姆斯·格雷戈里（James Gregory FRS，1638年11月－1675年10月）是一位苏格兰数学家和天文学家。他的姓氏有时拼作Gregorie，这是苏格兰的原始拼法。他提出了反射望远镜的早期实用设计——格里高利望远镜，并在三角学方面取得进展，发现了若干三角函数的无穷级数展开。

在他的著作《几何学的普遍部分》（Geometriae Pars Universalis，1668年）[1] 中，格雷戈里首次发表并证明了微积分基本定理（从几何角度陈述，并且仅限于后来版本定理所考虑曲线的一类特殊情况）。艾萨克·巴罗对此予以承认。[2][3][4][5][6][7]
\subsection{生平}
格雷戈里出生于1638年。他的母亲珍妮特是琼和大卫·安德森的女儿，父亲约翰·格雷戈里是一位苏格兰圣公会牧师。詹姆斯是他们三个孩子中最小的一个，出生在阿伯丁郡德鲁莫克的牧师住宅中，最初由母亲珍妮特·安德森（约1600–1668）在家中教育。正是他的母亲激发了他对几何学的兴趣，因为她的叔叔——亚历山大·安德森（Alexander Anderson, 1582–1619）曾是法国数学家维达的学生和著作编辑。在1651年父亲去世后，他的哥哥大卫接管了对他的教育。他先后就读于阿伯丁文法学校和马里斯哈尔学院（1653–1657），并于1657年获得艺术硕士学位（AM）。

1663年，他前往伦敦，结识了约翰·柯林斯以及苏格兰同胞罗伯特·穆雷（Robert Moray，皇家学会创始人之一）。1664年，他前往威尼斯共和国的帕多瓦大学，在途中途经佛兰德、巴黎和罗马。在帕多瓦，他与同乡、哲学教授詹姆斯·卡登黑德同住，并师从斯特凡诺·安杰利。他在帕多瓦发表的几何学论文《真正的圆和双曲线求积》（Vera circuli et hyperbolæ quadratura, 1667）和《几何学的普遍部分》（Geometriæ pars universalis, 1668），见证了安杰利及其他意大利著名数学家（如加布里埃莱·曼弗雷迪）对他的深刻影响[9]。

1668年返回伦敦后，他当选为皇家学会院士。同年晚些时候，他前往圣安德鲁斯大学，担任首任皇家数学教授，该职位由查理二世设立，很可能是应罗伯特·穆雷的请求。在圣安德鲁斯大学，他于1673年在实验室地板上铺设了第一条子午线，比格林尼治子午线确立早了200年，因此被认为“可以说圣安德鲁斯是时间的起点”[10][11]。

他先后在圣安德鲁斯大学和爱丁堡大学任教。

他与乔治·詹姆森（George Jameson，画家）的女儿、约翰·伯内特（John Burnet，来自阿伯丁埃尔里克）的遗孀玛丽结婚；他们的儿子詹姆斯成为阿伯丁国王学院的物理学教授。他是约翰·格雷戈里（John Gregory，1756年当选皇家学会院士）的祖父，大卫·格雷戈里（David Gregorie，1692年当选皇家学会院士）的叔叔，以及医生和发明家大卫·格雷戈里（1627–1720）的弟弟。

在担任爱丁堡大学数学讲座教授约一年后，詹姆斯·格雷戈里在与学生观测木星卫星时突发中风，几天后去世，享年36岁。
\subsection{出版著作}